\documentclass[12pt]{article}
\usepackage[T1]{fontenc}
\usepackage[utf8]{inputenc}
\usepackage{lmodern}
\usepackage{textcomp}
\usepackage{enumitem}
\usepackage{geometry}
\geometry{margin=1in}
\begin{document}
\begin{center}
Answer Key - Economics - K-12 Level Assessment 21 \
\small (Answers are ordered exactly as in the question paper.)
\end{center}

\section*{Multiple Choice}
\begin{enumerate}[label=\arabic*.]
\item \textbf{Answer:} D
\\ \textit{Options:} A) Lower wages relative to other nations.; B) Lower taxes on corporate profits relative to other nations.; C) A higher interest rate on financial assets relative to other nations.; D) A higher rate of inflation relative to other nations.
\\ \textit{Note:} A higher rate of inflation relative to other nations can lead to increased prices for domestic goods, making imports more attractive and exports less competitive, thereby worsening the trade deficit.
\item \textbf{Answer:} A
\\ \textit{Options:} A) Demand for white Ford minivans; B) Demand for automobiles; C) Demand for Ford automobiles; D) Demand for American-made automobiles
\\ \textit{Note:} The demand for white Ford minivans is likely to be more elastic because it is a specific brand and model, meaning consumers can easily switch to alternative vehicles if the price rises, indicating a greater sensitivity to price changes.
\item \textbf{Answer:} D
\\ \textit{Options:} A) The cash you receive from babysitting your neighbor's kids; B) The sale of illegal drugs; C) The sale of cucumbers to a pickle manufacturer; D) The sale of a pound of tomatoes at a supermarket
\\ \textit{Note:} The sale of a pound of tomatoes at a supermarket is counted in GDP because it represents a final transaction of goods produced within the economy, reflecting the total market value of all finished goods and services sold.
\item \textbf{Answer:} C
\\ \textit{Options:} A) Increasing the discount rate; B) Increasing the reserve ratio; C) Buying government securities; D) Lowering tariffs
\\ \textit{Note:} Buying government securities increases the money supply and lowers interest rates, which can lead to a depreciation of the U.S. dollar, making U.S. exports cheaper and more competitive in the global market.
\item \textbf{Answer:} A
\\ \textit{Options:} A) the value of the dollar will tend to appreciate.; B) the value of the dollar will tend to depreciate.; C) exchange rates will be affected but not the value of the dollar.; D) the exchange rate will not be affected.
\\ \textit{Note:} When interest rates rise in the United States, it attracts foreign investment seeking higher returns, increasing demand for the dollar and causing its value to appreciate relative to other currencies.
\end{enumerate}

\section*{True / False}
\begin{enumerate}[label=\arabic*.]
\setcounter{enumi}{5}
\item \textbf{Answer:} False
\\ \textit{Options:} A) True; B) False
\\ \textit{Note:} Deadweight loss can occur in various types of markets, including monopolistic or oligopolistic markets, when prices are above or below the equilibrium level, not just in perfectly competitive markets.
\item \textbf{Answer:} False
\\ \textit{Options:} A) True; B) False
\\ \textit{Note:} In capitalist market economies, private individuals and businesses own land and capital resources, while the government's role is primarily to regulate and facilitate the market rather than own all resources.
\item \textbf{Answer:} True
\\ \textit{Options:} A) True; B) False
\\ \textit{Note:} Higher taxes reduce consumer spending and aggregate demand, indicating contractionary fiscal policy, while buying Treasury securities by the central bank increases the money supply, indicating expansionary monetary policy, thus combining both effects as stated.
\item \textbf{Answer:} True
\\ \textit{Options:} A) True; B) False
\\ \textit{Note:} The statement is true because an increase in autonomous private investment leads to a multiplied effect on aggregate demand through the spending multiplier, which, in this case, assumes a multiplier of 10.
\item \textbf{Answer:} False
\\ \textit{Options:} A) True; B) False
\\ \textit{Note:} A kinked demand curve indicates that firms in an oligopoly face a downward-sloping demand curve above the kink, suggesting that price is typically greater than marginal cost (P \textgreater{} MC) in equilibrium, as firms are reluctant to change prices due to the expected reactions of competitors.
\end{enumerate}

\section*{Long Form Response}
\begin{enumerate}[label=\arabic*.]
\setcounter{enumi}{10}
\item \textbf{Answer:} Increased investment in capital infrastructure, such as roads, bridges, and public transportation, can significantly boost economic growth in the United States. When the government or private sector invests more in these long-lasting projects, it creates jobs, stimulates demand for materials, and enhances productivity by improving transportation and communication systems. At the same time, a reduced consumption of nondurable goods and services, like fast food or clothing, can free up more savings for investment. This means that people are spending less on items that do not last long and instead are contributing to projects that provide long-term benefits for the economy. Ultimately, this combination can lead to a stronger, more resilient economy as it focuses on sustainable growth and infrastructure development.
\\ \textit{Note:} correct because increased investment in capital infrastructure drives job creation and productivity improvements, while reduced consumption of nondurable goods allows for greater savings and investment potential, both of which contribute to faster economic growth.
\item \textbf{Answer:} The GDP deflator is a crucial measure used to assess changes in the price level of all goods and services produced in a country. It reflects the ratio of nominal GDP, which is measured at current prices, to real GDP, which is adjusted for inflation. By comparing these two figures, the GDP deflator indicates whether the overall price level in the economy has increased or decreased over time. This measure is significant because it helps economists and policymakers understand inflationary trends, allowing them to make informed decisions regarding monetary policy and economic planning. In essence, the GDP deflator serves as a vital tool for gauging the health of an economy and tracking its performance over time.
\\ \textit{Note:} correct because it accurately describes the GDP deflator as a measure that compares nominal GDP to real GDP, thereby providing insight into inflation and changes in the overall price level of national products.
\end{enumerate}

\vfill
\noindent\textit{End of Answer Key}
\end{document}